\documentclass[11pt,a4paper,reqno]{amsart}
\usepackage[T1]{fontenc}
\usepackage{lmodern}
\usepackage{amsmath,amssymb,amsthm,mathtools}
\usepackage{microtype}
\usepackage{hyperref}
\usepackage[left=1.12in,right=1.12in,top=0.98in,bottom=0.98in]{geometry}
\hypersetup{hidelinks,pdftitle={Cycles without clumps: a coreless etale correspondence in characteristic two},pdfauthor={OpenAI}}
\allowdisplaybreaks[2]
\setlength{\emergencystretch}{2em}
\setlength{\parskip}{0pt}
\newtheorem{theorem}{Theorem}[section]
\newtheorem{proposition}[theorem]{Proposition}
\newtheorem{corollary}[theorem]{Corollary}
\newtheorem{lemma}[theorem]{Lemma}
\newtheorem{remark}[theorem]{Remark}
\newtheorem{question}[theorem]{Question}
\newcommand{\kk}{k}
\newcommand{\FF}{\mathbf F}
\newcommand{\PP}{\mathbf P}
\newcommand{\Jac}{\operatorname{Jac}}
\newcommand{\Aut}{\operatorname{Aut}}
\newcommand{\Gal}{\operatorname{Gal}}
\newcommand{\Spec}{\operatorname{Spec}}
\newcommand{\Div}{\operatorname{div}}
\newcommand{\ord}{\operatorname{ord}}
\newcommand{\Pic}{\operatorname{Pic}}
\newcommand{\Tr}{\operatorname{Tr}}
\newcommand{\Frob}{\operatorname{Frob}}
\newcommand{\id}{\mathrm{id}}
\newcommand{\et}{\mathrm{\acute et}}
\newcommand{\GG}{\mathcal G}
\title[Cycles without clumps]{Cycles without clumps:\\ a coreless \texorpdfstring{$\acute{e}$}{e}tale correspondence\\ in characteristic two}
\author{OpenAI}
\date{20 November 2026 --- research preprint, version 0.2}
\subjclass[2020]{14H30, 14H10, 11G20, 05C25, 14G17}
\keywords{Finite \(\acute{e}\)tale correspondence, coreless correspondence, clump, generic graph, Hermitian curve, Cartier operator, \(p\)-rank}
\begin{document}
\begin{abstract}
We construct a symmetric finite \(\acute{e}\)tale self-correspondence of a
genus-two curve over $\overline{\FF}_2$ whose two maps have
degree $25$. The correspondence is minimal, has no core, and admits no
nonempty finite saturated set (clump). Its generic correspondence graph
nevertheless contains a $K_{5,5}$, and cannot be three-arc-transitive.
Thus one explicit construction gives negative answers to Questions
5.17, 6.11, and 9.7 in Krishnamoorthy's \emph{Correspondences without a
core}. The clump exclusion rests on a function-field incompatibility
between two wild quotient actions, which we formulate as a general
criterion. We also record a Cartier-theoretic $p$-rank obstruction and
produce further clump-free correspondences by abelian base change.
\end{abstract}
\maketitle

\section{Introduction and statements}
Let $k$ be an algebraically closed field. A \emph{correspondence of curves}
is a diagram
\begin{equation}\label{eq:correspondence}
  X\xleftarrow{\ f\ }N\xrightarrow{\ g\ }Y
\end{equation}
with finite separable maps of smooth projective connected curves. It is
\emph{\'{e}tale} when both maps are finite \(\acute{e}\)tale. It is
\emph{coreless} if
\[
  f^*k(X)\cap g^*k(Y)=k\qquad\text{inside }k(N).
\]
A \emph{clump} is a nonempty finite set $S\subseteq N(k)$ which is
saturated under both maps:
\[
   f^{-1}f(S)=S=g^{-1}g(S).
\]
A correspondence is \emph{minimal} when $k(N)$ is generated by its two
endpoint fields. Our conventions for the generic graph agree with
\cite[\S5]{Kri18}.

Three questions in \cite{Kri18} ask whether the generic graph of a
minimal coreless correspondence must be acyclic (Question 5.17), whether
the canonical action on the generic graph of a minimal symmetric coreless
\(\acute{e}\)tale self-correspondence is infinity-arc-transitive
(Question 6.11), and whether a projective coreless
\(\acute{e}\)tale correspondence in characteristic $p>0$ must have a
clump (Question 9.7).  Our construction gives counterexamples to all three questions
using the same source curve.  The first two consequences follow from
the finite subgroup geometry; the third requires a separate argument
with wild quotient groups.

\begin{theorem}[A single symmetric example]\label{thm:main}
Set $k=\overline{\FF}_2$. There is a smooth projective curve $N$ of
genus $26$ and two finite \(\acute{e}\)tale maps of degree $25$ to
\[
     X:\qquad y^2+y=x^5,\qquad g(X)=2,
\]
with the following properties:
\begin{enumerate}
\item The correspondence $X\xleftarrow{f}N\xrightarrow{g}X$ is symmetric,
minimal and coreless.
\item It has no nonempty clump.
\item Its generic graph contains a complete bipartite subgraph $K_{5,5}$.
The subgraph has $25$ edges, $100$ unoriented four-cycles, and cycle rank
$16$.
\item The action on its infinite connected generic graph is not
three-arc-transitive (in fact, no graph automorphism group can act
three-arc-transitively).
\end{enumerate}
\end{theorem}

The essential distinction is between a \emph{finite graph of incidences}
and a \emph{finite subset closed under all incidence operations}.  The
$K_{5,5}$ of Theorem~\ref{thm:main} is not a clump: each of its vertices
has five edges within the displayed graph but valency $25$ in the full
generic graph.

The construction is explicit. Let
\begin{equation}\label{eq:hermitian}
 H:\quad U^5+V^5=W^5
\end{equation}
be the Fermat--Hermitian plane quintic, a curve of genus six.  There are
\(\acute{e}\)tale degree-five quotient maps $\pi,\psi:H\to X$ and an
involution $\gamma$ of $X$ such that
\begin{equation}\label{eq:fiberprod}
 N=H\times_{\pi,X,\gamma\pi}H,
 \qquad f=\psi\circ p_1,\quad g=\psi\circ p_2.
\end{equation}
The source is connected, and exchanging the two factors realizes the
symmetry.  The proof that it has no clump is the main content of the paper.

A second, characteristic-$p$ statement gives an independent diagnostic.
Here $\gamma_p(C)$ denotes the $p$-rank of a curve over
$\overline{\FF}_p$.

\begin{proposition}[Cartier intersection criterion]\label{prop:rank}
For any finite \(\acute{e}\)tale correspondence
$X\xleftarrow{f}Z\xrightarrow{g}Y$ over $\overline{\FF}_p$, with
$g(X),g(Y)\ge 2$, put
\[
 r=\gamma_p(X)+\gamma_p(Y)-\gamma_p(Z).
\]
If $r\ge 1$, the correspondence has a nonempty clump. If $r\ge2$,
it has a core. In particular,
\begin{align}
 \text{no clump}&\ \Longrightarrow\ 
 \gamma_p(Z)\ge\gamma_p(X)+\gamma_p(Y),\label{eq:rank-ineq}\\
 \text{coreless}&\ \Longrightarrow\ 
 \gamma_p(Z)\ge\gamma_p(X)+\gamma_p(Y)-1.\label{eq:core-ineq}
\end{align}
\end{proposition}

This proposition is included as a short structural criterion, not as an
assertion of bibliographic priority. Its proof is given in
\S\ref{sec:cartier}.  Our example has $\gamma_2(X)=0$; thus the
criterion is compatible with its absence of clumps.

\section{The geometric construction}\label{sec:construction}
Fix $\omega\in\FF_4$ with $\omega^2+\omega=1$ and a primitive fifth
root of unity $\zeta\in\FF_{16}$. Write $u=U/W$, $v=V/W$ on $H$ and set
\begin{equation}\label{eq:pi}
  \pi:H\longrightarrow X,\qquad x=uv,\quad y=u^5.
\end{equation}
Indeed $u^5+v^5=1$ implies
\[
   y^2+y=u^5(u^5+1)=u^5v^5=x^5.
\]
The automorphism
\[
 d(U:V:W)=(\zeta U:\zeta^{-1}V:W)
\]
has order five, fixes the functions in~\eqref{eq:pi} and acts freely
on $H$: the only fixed points of any nonidentity power are coordinate
vertices of $\PP^2$, none of which lies on the Fermat quintic. Thus
$\pi$ is a cyclic \(\acute{e}\)tale cover of degree five.

Define
\begin{equation}\label{eq:gamma}
  \gamma(x,y)=(x+1,\,y+x^2+x+\omega).
\end{equation}
Since $(x+1)^5+x^5=x^4+x+1$ and
$(x^2+x+\omega)^2+(x^2+x+\omega)=x^4+x+1$,
this is an involution of $X$. It fixes the unique point $O$ at infinity.

For the second quotient, take $\FF_{16}=\FF_2[\theta]/(\theta^4+\theta+1)$
and $\mathbf v=(1,\theta,\theta^2)$. For the Hermitian pairing
\[
  \langle z,w\rangle=\sum_{i=1}^3z_iw_i^4,
\]
we have $\langle\mathbf v,\mathbf v\rangle=0$, hence
\[
  T(z)=z+\langle z,\mathbf v\rangle\mathbf v
\]
is an involutory unitary transvection.  Put
$F=\langle d\rangle$, $F'=TFT^{-1}$, and
\[
    \psi=\pi\circ T^{-1}:H\longrightarrow X.
\]
Since $T$ is an automorphism of $H$, the map $\psi$ is another cyclic
\(\acute{e}\)tale quotient of degree five.  Equation~\eqref{eq:fiberprod}
now defines the desired symmetric correspondence.

\begin{lemma}[Connectedness]\label{lem:connected}
The fibre product $N$ in~\eqref{eq:fiberprod} is connected.  It is a
$C_5\times C_5$-Galois \(\acute{e}\)tale cover of $X$ of degree
$25$, and $g(N)=26$. Both exterior maps $f,g$ have degree $25$.
\end{lemma}
\begin{proof}
The Kummer equation for $H/X$ is $u^5=y$.  Let $P_0=(0,0)$;
then $\Div(y)=5(P_0-O)$. The divisor class
$\xi=[P_0-O]\in\Jac(X)[5]$ has exact order five: it is not
trivial, since a degree-one function cannot exist on $X$.
The conjugate cover has Kummer class $\gamma^*\xi$. Let
$Q=\gamma(P_0)=(1,\omega)$. If the two cyclic Kummer lines
coincided, $[Q-O]=j[P_0-O]$ for some
$j\in\{\pm1,\pm2\}$.
For $j=\pm1$, the injectivity of the degree-one Abel map gives
$Q=P_0$ or $Q=\iota P_0=(0,1)$, both impossible, where
$\iota(x,y)=(x,y+1)$ is the hyperelliptic involution.
For $j=\pm2$, replace $P_0$ if needed by $\iota P_0$ and obtain
$Q+O\sim2R$ with $R\in\{P_0,\iota P_0\}$.  The two effective
degree-two divisors are distinct. By the genus-two Riemann--Roch
theorem the linear system must be the hyperelliptic pencil, so
$2R\sim2O$, contrary to the order-five property of $[R-O]$.
Thus the Kummer lines are distinct, and the corresponding cyclic
extensions are linearly disjoint.  The fibre product is connected and
has group $C_5^2$ over $X$. Since the cover is \(\acute{e}\)tale,
Riemann--Hurwitz gives $g(N)=1+25(g(X)-1)=26$.
The two projections $p_i:N\to H$ have degree five; composing with
$\psi$ gives the asserted exterior degrees.
\end{proof}

\section{Two wild quotient groups}\label{sec:wild}
The clump exclusion rests on the interaction between two finite
translation groups. We record the exact finite group input rather than
appealing to a classification of unitary groups.

\begin{lemma}[Finite group certificate]\label{lem:certificate}
Let $G=\langle F,F'\rangle\subset\Aut(H)$, with $F,F'$ as in
\S\ref{sec:construction}. Then $G$ contains
\[
 \Lambda:(U:V:W)\longmapsto(\zeta U:V:W)
\]
and the following group $P_H$ of order $64$. On the affine chart
\[
  s=\frac{W}{U+V},\qquad t=\frac{U}{U+V},
  \qquad t^4+t=s^5+1,
\]
its elements act as
\begin{equation}\label{eq:PH}
 (s,t)\longmapsto
 (s+a,\ t+a^4s+b),\qquad
 a,b\in\FF_{16},\quad b^4+b=a^5.
\end{equation}
The group containments admit the short explicit word certificate below.
\end{lemma}
\begin{proof}
The formulas in~\eqref{eq:PH} preserve $t^4+t=s^5+1$:
in characteristic two,
\[
 (s+a)^5=s^5+a s^4+a^4s+a^5.
\]
For each $a$ there are four admissible $b$, giving $64$ distinct
transformations. Let $e=TdT^{-1}$ and abbreviate $D=d^{-1}$,
$E=e^{-1}$. A word is evaluated as a matrix product from left
to right, modulo scalar matrices. In homogeneous coordinates the
element in~\eqref{eq:PH} is represented by
\[
 M(a,b)=\begin{pmatrix}
  1+b&b&a^4\\ b&1+b&a^4\\ a&a&1
 \end{pmatrix}.
\]
The following five identities hold over $\FF_{16}$, with
$\theta^4+\theta+1=0$ and $\zeta=\theta^3$:
\begin{center}
\begin{tabular}{@{}ll@{}}
  element & word in $d,e,D,E$\\\hline
  $M(1,\theta^2)$ & \texttt{eeDeeddEDD}\\
  $M(\theta,\theta^3+\theta)$ & \texttt{DeDEDDeded}\\
  $M(\theta^2,\theta^3)$ & \texttt{edeDDee}\\
  $M(\theta^3,\theta^2)$ & \texttt{eddEddeed}\\
  $\Lambda=\operatorname{diag}(\zeta,1,1)$ & \texttt{EdeDDedEDD}
\end{tabular}
\end{center}
These identities are checkable by direct $3\times3$ matrix
multiplication; a second, independent finite-field verification
is provided by the accompanying script.

The first four elements generate all of $P_H$:
their $s$-translation parameters $1,\theta,\theta^2,\theta^3$
form an $\FF_2$-basis of $\FF_{16}$.
The commutator of $M(a,b)$ and $M(c,d)$ has parameter
$(0,a^4c+c^4a)$.  In particular, the commutators for
$(a,c)=(1,\theta)$ and $(1,\theta^3)$ yield the two
independent central translations $1$ and
$\theta^2+\theta+1$ in $\FF_4$.
Thus the subgroup they generate has at least $16\cdot4=64$
elements, and it is contained in $P_H$ of order $64$.
The last word gives $\Lambda\in G$. This proves the required
containments without enumerating the full generated group.

\end{proof}

\begin{lemma}[Fixed fields]\label{lem:fixed-fields}
The invariant fields of the two groups relevant below are
\[
  k(H)^{P_H}=k(B),\quad B=s^{16}+s,
  \qquad
  k(X)^{P_X}=k(A),\quad A=x^{16}+x,
\]
where $\lambda:(x,y)\mapsto(\zeta x,y)$ and $P_X\subset\Aut(X)$
is the subgroup generated by the five conjugates
$\lambda^i\gamma\lambda^{-i}$ for $0\le i<5$. Its order is $32$.
Moreover,
\begin{equation}\label{eq:diff-identity}
    dx=s^2\,ds,\qquad dA=dx,\quad dB=ds
\end{equation}
on $H$.
\end{lemma}
\begin{proof}
The group $P_H$ fixes $B$ because $a^{16}=a$.  The affine equation
$t^4+t=s^5+1$ gives $[k(H):k(s)]=4$ and
$[k(s):k(B)]=16$. Hence $[k(H):k(B)]=64$, the group order.

The conjugates $\lambda^i\gamma\lambda^{-i}$ translate $x$ by
$\zeta^i$. Since $1,\zeta,\zeta^2,\zeta^3$ span $\FF_{16}$ over
$\FF_2$, the subgroup they generate surjects onto the additive
translations of $\FF_{16}$. Its kernel lies in the hyperelliptic
involution $\iota:(x,y)\mapsto(x,y+1)$. A commutator of the
translations by $1$ and $\zeta$ is precisely $\iota$, since the
constant added to $y$ is
$\zeta^2+\zeta+\zeta^{-2}+\zeta^{-1}=1$.
It follows that $|P_X|=32$. Since $A$ is invariant and
$[k(X):k(A)]=2\cdot16=32$, the fixed field is $k(A)$.

Finally, $x=(t^2+t)/s^2$ on $H$. Differentiating
$t^4+t=s^5+1$ gives $dt=s^4ds$, so
$dx=s^{-2}dt=s^2ds$. The remaining identities follow from
$d(s^{16}+s)=ds$ and $d(x^{16}+x)=dx$.
\end{proof}

\section{A function-field obstruction to clumps}\label{sec:noclump}
We require two elementary intersection lemmas.

\begin{lemma}\label{lem:int1}
Inside $k(H)$,
\begin{equation}\label{eq:int1}
       k(x)\cap k(s)=k(E),\qquad E=x^6+x.
\end{equation}
\end{lemma}
\begin{proof}
Put $w=t^2+t$ and $D=s^5+1$, so $w^2+w=D$.
A direct reduction yields
\begin{equation}\label{eq:E}
 E=\frac{D+D^2+D^3}{s^{12}}\in k(s).
\end{equation}
The extension $k(H)/k(x)$ has degree ten. Its full automorphism
group is the dihedral group $D_{10}$ generated by $F$ and
$U\leftrightarrow V$. Similarly, $k(H)/k(s)$ is the Galois extension
of degree four generated by $t\mapsto t+c$ for $c\in\FF_4$.
Let $J$ be the subgroup of $\Aut(H)$ generated by these two Galois
groups. Both groups fix $E$, so
$|J|\mid[k(H):k(E)]=60$. Since their orders are $10$ and $4$,
$|J|$ is $20$ or $60$.

If $|J|=20$, its subgroup $D_{10}$ has index two, hence is normal.
Its unique characteristic subgroup of order five, $F$, would also be
normal in $J$. But the element $t\mapsto t+c$, with
$c\in\FF_4\setminus\FF_2$, is represented on $(U,V,W)$ by a matrix
with upper-left block
\[
  \begin{pmatrix}1+c&c\\c&1+c\end{pmatrix},
\]
which does not normalize $F$ (it fails to permute the eigenlines of
$d$). This is a contradiction. Thus $|J|=60$.
The fixed field of $J$ is $k(E)$ by the equality of extension degrees,
which is the intersection of the two individual fixed fields.
\end{proof}

\begin{lemma}\label{lem:int2}
In the rational function field $k(x)$,
\begin{equation}\label{eq:int2}
   k(x^{16}+x)\cap k(x^6+x)=k.
\end{equation}
\end{lemma}
\begin{proof}
Write $A=x^{16}+x$ and $E=x^6+x$, so $A'=E'=1$.
If a nonconstant element lies in both rational subfields, express it
in reduced rational form as
$P(A)/Q(A)=R(E)/S(E)$.  Coprimeness remains valid after polynomial
substitution, by B\'ezout identities. Comparing numerators and
denominators in $k[x]$ therefore produces a common nonconstant
polynomial, say $P_0(A)=R_0(E)$. Choose one of least positive degree
(as a polynomial in $x$).

Differentiating and using $A'=E'=1$ gives
$P_0'(A)=R_0'(E)$. By minimality, this common polynomial is constant,
say $c$. If $c=0$, both $P_0,R_0$ are squares over the algebraically
closed field $k$, and taking their square roots yields a common
polynomial of smaller positive degree. Thus $c\ne0$.
Write
\[
 P_0(T)=cT+P_1(T)^2,\qquad
 R_0(T)=cT+R_1(T)^2.
\]
Since $A+E=(x^8+x^3)^2$, taking square roots gives
\begin{equation}\label{eq:sqrt1}
 P_1(A)+R_1(E)=\sqrt c\,(x^8+x^3).
\end{equation}
Differentiate again. In characteristic two the derivative of any
polynomial is a square polynomial; hence taking square roots of the
result yields polynomials $P_2,R_2$ satisfying
\begin{equation}\label{eq:sqrt2}
 P_2(A)+R_2(E)=c^{1/4}x.
\end{equation}
The left side has the same value at $x=0$ and $x=1$, since
$A(0)=A(1)=E(0)=E(1)=0$; the right side does not. Contradiction.
\end{proof}

\begin{lemma}[Incompatible wild quotients]\label{lem:wild-obstruction}
Let $k=\overline{\FF}_p$ and let $\pi:H\to X$ be a finite
\(\acute{e}\)tale map with $g(X)\ge2$. Suppose finite $p$-subgroups
$P_X\le\Aut(X)$ and $P_H\le\Aut(H)$ have rational invariant
fields $k(X)^{P_X}=k(u)$ and $k(H)^{P_H}=k(v)$, where
$du,dv\ne0$. Write $\pi^*du=\rho\,dv$, with $\rho\in k(H)^\times$.
Assume
\begin{equation}\label{eq:wild-criterion}
 k(u)\cap k(v,\rho)=k,
 \qquad \rho^n\notin k(v)\quad\text{for every }n\ge1.
\end{equation}
Then no nonzero effective divisor $D$ on $X$ can be
$P_X$-invariant while $\pi^*D$ is $P_H$-invariant.
\end{lemma}
\begin{proof}
Write $r=\deg D>0$ and $g=g(X)$.  Since every point of
$\Pic^0(X)(k)$ is torsion, choose $e>0$ such that
\[
 e\bigl((2g-2)D-rK_X\bigr)\sim0.
\]
Put $a=e(2g-2)$ and $n=er$. There is a rational
$n$-fold canonical differential $\eta_X$ with
$\Div(\eta_X)=aD$. The group $P_X$ preserves its divisor,
so it acts on the line $k\eta_X$ through a character
$P_X\to k^\times$. This character is trivial in characteristic
$p$, because $P_X$ is a $p$-group. The invariant differential
$du$ therefore gives
$\eta_X=h(u)(du)^n$ for some $h\in k(u)^\times$.

The pullback $\pi^*\eta_X$ has divisor $a\pi^*D$ and is
similarly $P_H$-invariant. Thus
$\pi^*\eta_X=q(v)(dv)^n$ for some $q\in k(v)^\times$.
Comparing expressions yields
\[
          h(u)\rho^n=q(v).
\]
Consequently $h(u)\in k(u)\cap k(v,\rho)=k$, which forces
$\rho^n\in k(v)$, contradicting~\eqref{eq:wild-criterion}.
\end{proof}

\begin{theorem}[No finite saturated subset]\label{thm:noclump}
The correspondence in~\eqref{eq:fiberprod} has no nonempty clump.
In particular it is coreless.
\end{theorem}
\begin{proof}
Suppose $S\subset N(k)$ is a nonempty clump. The interchange of the
two factors of $N$ interchanges the exterior maps $f,g$; replacing $S$
by its union with its image under this interchange yields another
nonempty clump which is invariant under interchange.
Since the exterior legs factor through $p_1,p_2$, the set $S$ is
saturated for both coordinate projections. Its two images in $H$
are the same reduced finite set $D_H$, by interchange invariance.
The coordinate deck groups together generate
$\Gal(N/X)\cong C_5^2$; thus there exists a nonempty reduced divisor
$D_X$ on $X$ such that
\[
  S=(\pi p_1)^{-1}(D_X),\qquad
  D_H=\pi^{-1}(D_X),\qquad \gamma(D_X)=D_X.
\]
The equality $f^{-1}f(S)=S$ implies that $D_H$ is also saturated
for the quotient $\psi:H\to H/F'\cong X$. Since $D_H$ is already
$F$-stable, it is stable under $G=\langle F,F'\rangle$.
Lemma~\ref{lem:certificate} makes it $P_H$-stable and gives
$\Lambda$-stability. The latter descends to
$\lambda:(x,y)\mapsto(\zeta x,y)$; together with $\gamma$ it
makes $D_X$ stable under $P_X$.

Apply Lemma~\ref{lem:wild-obstruction} to
\[
 u=A=x^{16}+x,\qquad v=B=s^{16}+s,\qquad \rho=s^2,
\]
using~\eqref{eq:diff-identity}. The two field-intersection
lemmas show
\[
 k(A)\cap k(B,s^2)=k(A)\cap k(s)=k;
\]
indeed $s=B+(s^2)^8$. Furthermore, $s^{2n}\notin k(B)$ for
any $n>0$: the translation $s\mapsto s+1$ fixes $B$ but does
not fix $s^{2n}$, as their zero divisors have different support.
Thus the wild-quotient criterion contradicts the existence of
$D_X$. No nonempty clump exists.

If a nonconstant common function lay in
$f^*k(X)\cap g^*k(X)$, the support of one of its nonempty finite
geometric fibres would be a nonempty set saturated under both legs.
The result just proved excludes this. Hence the correspondence is
coreless.
\end{proof}

\begin{corollary}[No invariant pluricanonical forms]
\label{cor:pluri}
For the correspondence of Theorem~\ref{thm:main} and every
$m\ge1$, there are no nonzero sections
$\alpha,\beta\in H^0(X,(\Omega_X^1)^{\otimes m})$ with
$f^*\alpha=g^*\beta$.
\end{corollary}
\begin{proof}
Any nonzero such common section would have a nonempty finite
zero divisor, whose support is saturated for both maps by
\(\acute{e}\)taleness, contradicting
Theorem~\ref{thm:noclump}.
\end{proof}

\section{Minimality and cycles in the generic graph}\label{sec:graph}
We first make explicit why the preceding example satisfies the
minimality hypothesis of the graph questions.

\begin{proposition}\label{prop:minimal}
The correspondence in Theorem~\ref{thm:main} is minimal.
\end{proposition}
\begin{proof}
Let $K=k(N)$ and let $K_1,K_2\subset K$ be the two exterior
endpoint fields. Put $M=K_1K_2$. The degree
$d=[M:K_1]=[M:K_2]$ divides $25$.
Each tower $K/K_i$ is a composite of two cyclic extensions of
order five. Hence the Galois closure of $K/K_i$, and therefore that
of $M/K_i$, has a $5$-group as Galois group.

If $d=1$, then $K_1=K_2$, contradicting corelessness.
If $d=5$, a transitive $5$-subgroup of $S_5$ is necessarily $C_5$.
Both $M/K_i$ would therefore be cyclic Galois extensions of degree
five. Since $M/K_i$ is \(\acute{e}\)tale of degree five over a
genus-two curve, the smooth projective model of $M$ has genus six.
Its automorphism group over $k$ is finite. The subgroup generated by
$\Gal(M/K_1)$ and $\Gal(M/K_2)$ would therefore be finite, and its
fixed field would be a function field of transcendence degree one
contained in $K_1\cap K_2$. This is impossible.
Thus $d=25$ and $M=K$.
\end{proof}

\begin{proposition}[Finite deck groups force a generic cycle]
\label{prop:cycle}
Let $X\xleftarrow{f}Z\xrightarrow{g}Y$ be a minimal finite
separable correspondence over an algebraically closed field, with
$g(Z)\ge2$. Write
\[
\begin{gathered}
 K=k(Z),\quad L=f^*k(X),\quad R=g^*k(Y),\\
 A=\Aut(K/L),\quad B=\Aut(K/R),\quad \Gamma=\langle A,B\rangle.
\end{gathered}
\]
If $|A|,|B|\ge2$, the generic graph contains a finite connected
bipartite subgraph of cycle rank
\begin{equation}\label{eq:cycle-rank}
    1+|\Gamma|\left(1-\frac1{|A|}-\frac1{|B|}\right)\ge1.
\end{equation}
\end{proposition}
\begin{proof}
Since $K=LR$, we have $A\cap B=\{1\}$. The group
$\Gamma$ is finite because the automorphism group of a hyperbolic
smooth projective curve is finite.
Consider the edges of the full generic graph given by the field
embeddings $\sigma:K\hookrightarrow\overline{k(K)}$ for
$\sigma\in\Gamma$. Two such edges share their left vertex exactly
when their restrictions to $L$ agree, equivalently when they belong to
the same right coset of $A$; similarly, right vertices are right
cosets of $B$.  Minimality makes these edges distinct.
The resulting finite coset-incidence graph has $|\Gamma|$ edges and
$|\Gamma|/|A|+|\Gamma|/|B|$ vertices; it is connected because $A,B$
generate $\Gamma$. Each vertex has valency $|A|$ or $|B|$, both at
least two. The Euler formula $b_1=E-V+1$ proves the displayed
cycle-rank formula. These edges lie in the distinguished connected
component of the generic graph, since each element of $\Gamma$ is
obtained by alternating generators in $A$ and $B$.
\end{proof}

\begin{proof}[Proof of Theorem~\ref{thm:main}, graph assertions]
The two coordinate deck groups $A_0=\Gal(N/H)$ (for $p_1$) and
$B_0=\Gal(N/H)$ (for $p_2$) are commuting groups of order five,
with trivial intersection. They lie in the two exterior deck groups,
respectively. Their product is $C_5^2$.
Minimality implies that distinct pairs of elements of $A_0\times
B_0$ give distinct graph edges, with the blue endpoint depending only
on the $B_0$-coordinate and the red endpoint only on the
$A_0$-coordinate. We obtain all $25$ edges of a $K_{5,5}$.
It has $\binom52^2=100$ four-cycles and
$b_1=25-10+1=16$.

By \cite[Proposition 5.10]{Kri18}, a coreless correspondence has an
infinite generic graph. This graph is connected and locally finite,
with valency $25$ at every vertex. Its four-cycles contain
non-backtracking three-arcs whose endpoints are adjacent.
An infinite locally finite connected graph also has a geodesic
three-arc whose endpoints are at distance three. Adjacency of the
endpoints is preserved under any graph automorphism, so the graph
cannot be three-arc-transitive. The involution interchanging the two
coordinates of $N$ exhibits symmetry of the correspondence.
\end{proof}

\section{The Cartier rank obstruction}\label{sec:cartier}
We prove Proposition~\ref{prop:rank} for arbitrary characteristic
$p>0$.

\begin{proof}[Proof of Proposition~\ref{prop:rank}]
Let $C_C$ denote the Cartier operator on
$H^0(C,\Omega_C^1)$, and let
\[
  W_C=\bigcap_{m\ge0}\operatorname{im}(C_C^m).
\]
Over the algebraically closed ground field this is a vector space of
dimension $\gamma_p(C)$. Cartier commutes with the pullback of a
finite \(\acute{e}\)tale morphism, and its restriction to $W_C$ is
bijective. Thus $f^*W_X$ and $g^*W_Y$ are subspaces of $W_Z$ and
\[
 \dim(f^*W_X\cap g^*W_Y)
 \ge \gamma_p(X)+\gamma_p(Y)-\gamma_p(Z)=r.
\]
If $r\ge1$, choose a nonzero
$\omega=f^*\omega_X=g^*\omega_Y$ in the intersection.
A nonzero holomorphic differential on a curve of genus at least two
has a nonempty finite zero divisor. Because $f,g$ are
\(\acute{e}\)tale,
\[
 \operatorname{Supp}\Div(\omega)
 =f^{-1}\bigl(\operatorname{Supp}\Div(\omega_X)\bigr)
 =g^{-1}\bigl(\operatorname{Supp}\Div(\omega_Y)\bigr).
\]
This is a nonempty clump.

If $r\ge2$, choose independent $\omega_1,\omega_2$ in the
intersection. Their ratio is a nonconstant rational function on $Z$.
Because each $\omega_i$ is independently pulled back from both ends,
$\omega_1/\omega_2$ belongs to both endpoint function fields.
Thus the correspondence has a core. The contrapositives give
\eqref{eq:rank-ineq} and \eqref{eq:core-ineq}.
\end{proof}

\begin{remark}
The rank inequalities are necessary conditions, not equivalences.
They do not claim that a sufficiently large $p$-rank creates a
clump-free or coreless correspondence. They are intentionally
separated from the construction of Theorem~\ref{thm:main}; no
priority is asserted for this elementary dimension argument.
\end{remark}

\section{Propagation by abelian base change}\label{sec:propagation}
We record a concise extension which does not require further
Hermitian-group calculations.

\begin{proposition}\label{prop:propagation}
Let $k=\overline{\FF}_p$, and suppose that a connected finite
\(\acute{e}\)tale correspondence
$X\xleftarrow{f}N\xrightarrow{g}X$ has no core and no nonempty
clump. Then for every positive integer $n$ prime to $p$ and to the order of
the finite invariant Picard group scheme $\mathcal K$ defined in the proof, any
connected cyclic degree-$n$ \(\acute{e}\)tale cover $X_n\to X$
induces a connected finite \(\acute{e}\)tale correspondence
\[
 X_n\ \longleftarrow\ N_n\ \longrightarrow\ X_n,
 \qquad
 N_n=X_n\times_{X,f}N\times_{X,g}X_n,
\]
which has no nonempty clump and no core.  If the original
correspondence is symmetric, so is this one. For the example of
Theorem~\ref{thm:main},
\[
 g(X_n)=n+1,\quad g(N_n)=25n^2+1,
 \quad \deg(N_n\to X_n)=25n.
\]
In particular the construction furnishes such correspondences
in infinitely many genera.
\end{proposition}
\begin{proof}
By \cite[Lemma 8.9]{Kri18}, the group scheme
\[
 \mathcal K=\ker\left(\Jac(X)\times\Jac(X)
 \xrightarrow{\ f^*-g^*\ }\Jac(N)\right)
\]
is finite. Choose $n$ prime to $p$ and to the order of
$\mathcal K$.  Its $n$-torsion is trivial, so the pullback maps for
the two legs give an injection
\[
 H^1_{\et}(X,\mathbf Z/n)\oplus H^1_{\et}(X,\mathbf Z/n)
 \hookrightarrow H^1_{\et}(N,\mathbf Z/n).
\]
Here one identifies $\mu_n$ with the constant $\mathbf Z/n$
after choosing a primitive root of unity, and uses Kummer theory.
For a connected cyclic degree-$n$ cover, its character has exact
order $n$. The two pullback characters on $N$ are jointly
surjective onto $(\mathbf Z/n)^2$ by the displayed injection, so
$N_n$ is connected.

If $N_n$ had a nonempty clump, its projection to $N$ would give a
nonempty clump there. Indeed, the full fibre product contains every
pair of lifts over the two endpoints, and saturation on either side
persists after projection. This contradicts the hypothesis. Absence
of clumps also excludes a core.
Riemann--Hurwitz gives the indicated genera and degrees in the
special example. The coordinate-exchange involution extends by
interchanging the two copies of $X_n$ when the seed is symmetric.
\end{proof}

\begin{corollary}[Two independently varying endpoints]
\label{cor:mixed}
In the setting of Proposition~\ref{prop:propagation}, let
$X_n\to X$ and $X_m\to X$ be connected cyclic
\(\acute{e}\)tale covers, with $nm$ prime to $p|\mathcal K|$.
Then the full fibre product
\[
 N_{n,m}=X_n\times_{X,f}N\times_{X,g}X_m
\]
is connected and defines a clump-free, coreless
\(\acute{e}\)tale correspondence between $X_n$ and $X_m$.
For the example of Theorem~\ref{thm:main}, the endpoint
and source genera are $n+1,m+1,25nm+1$, while the
respective leg degrees are $25m,25n$.
In particular, taking $n\ne m$ gives correspondences between
curves of different genera.
\end{corollary}
\begin{proof}
Apply the injection in Proposition~\ref{prop:propagation}
with $\ell=\operatorname{lcm}(n,m)$.  If the joint monodromy
$\pi_1(N)\to(\mathbf Z/n)\times(\mathbf Z/m)$ were not
surjective, a nontrivial character of its cokernel would
contradict this injection.  The remaining claims follow from
full fibre-product saturation and Riemann--Hurwitz as before.
\end{proof}

\section*{Proof provenance and verification}
The finite group containments follow from the five short matrix words
in Lemma~\ref{lem:certificate}; the accompanying exact Python code
checks them independently and also enumerates the ambient group. The no-clump argument and field-intersection
lemmas are written in full above and do not depend on bounded-degree
searches. The computer certificate verifies the specified subgroup
containments, not the entire geometric proof. This is an unrefereed
research preprint; an independent mathematical audit and a complete
bibliographic priority check remain necessary before journal
submission.

\begin{thebibliography}{9}
\bibitem{Kri18}
R.~Krishnamoorthy, \emph{Correspondences without a core},
Algebra \& Number Theory \textbf{12} (2018), no.~5, 1173--1214.
\href{https://doi.org/10.2140/ant.2018.12.1173}{doi:10.2140/ant.2018.12.1173}.
\bibitem{Serre}
J.-P.~Serre, \emph{Trees}, Springer Monographs in Mathematics,
Springer, Berlin, 2003.
\bibitem{Stichtenoth}
H.~Stichtenoth, \emph{Algebraic Function Fields and Codes},
2nd ed., Graduate Texts in Mathematics 254, Springer, 2009.
\end{thebibliography}
\end{document}
